\documentclass[aspectratio=169,10pt,fontset=fandol]{ctexbeamer}
\usepackage{amsmath,amssymb,mathtools,bm,mathrsfs}
\usepackage{graphicx,booktabs,array,tabularx,multirow}
\usepackage{tikz,pgfplots,xparse,xstring,mdframed}
\usetikzlibrary{arrows.meta,calc,angles,positioning,shapes.geometric,graphs}
\pgfplotsset{compat=1.18}

\setbeamerfont{normal text}{size=\normalsize}
\setbeamerfont{title}{size=\large}
\setbeamerfont{subtitle}{size=\normalsize}
\setbeamerfont{frametitle}{size=\large}
\setbeamerfont{framesubtitle}{size=\small}
\setbeamerfont{block title}{size=\normalsize}
\setbeamerfont{block body}{size=\normalsize}
\setbeamerfont{section in toc}{size=\normalsize}
\setbeamerfont{section in head/foot}{size=\footnotesize}
\setbeamerfont{author in head/foot}{size=\footnotesize}
\setbeamerfont{title in head/foot}{size=\footnotesize}
\setbeamerfont{date in head/foot}{size=\footnotesize}
\setbeamertemplate{navigation symbols}{}
\setbeamertemplate{itemize item}{\normalsize\textbullet}
\setbeamertemplate{itemize subitem}{\small\textbullet}
\setbeamercolor{structure}{fg=blue!65!black}
\setbeamercolor{block title}{fg=blue!65!black,bg=blue!8}
\setbeamercolor{block body}{fg=black,bg=blue!3}
\setbeamertemplate{footline}{\leavevmode\hbox to\paperwidth{\hfill{\footnotesize\insertframenumber}\hspace*{1em}}\vskip3pt}

% ctex fontset=fandol; system fonts are not redistributed.

\newenvironment{proofdetail}{\par\noindent\textbf{证明：}\quad}{\par}
\renewenvironment{solution}{\par\noindent\textbf{解：}\quad}{\par}
\title{布局回归}
\begin{document}
\section{归纳证明}
% FRAME-ID F001
\begin{frame}[t,fragile]{目录}
\normalsize
{\color{blue}\normalsize 1. 归纳证明\par}\vspace{0.22em}
\end{frame}
% FRAME-ID F002
\begin{frame}[t,fragile]{归纳证明}
\footnotesize
{\footnotesize
\textbf{证明：}
% SOURCE-BEGIN U0001 0 222
第1步，这里保留完整的推导条件及数学论证。这里保留完整的推导条件及数学论证。这里保留完整的推导条件及数学论证。\[a_{1}=a_{0}+1\]

第2步，这里保留完整的推导条件及数学论证。这里保留完整的推导条件及数学论证。这里保留完整的推导条件及数学论证。\[a_{2}=a_{1}+1\]

第3步，这里保留完整的推导条件及数学论证。这里保留完整的推导条件及数学论证。这里保留完整的推导条件及数学论证。\[a_{3}=a_{2}+1\]


% SOURCE-END U0001
}
\end{frame}
% FRAME-ID F003
\begin{frame}[t,fragile]{归纳证明（续1）}
\footnotesize
{\footnotesize
\textbf{证明：}
% SOURCE-BEGIN U0001 222 444
第4步，这里保留完整的推导条件及数学论证。这里保留完整的推导条件及数学论证。这里保留完整的推导条件及数学论证。\[a_{4}=a_{3}+1\]

第5步，这里保留完整的推导条件及数学论证。这里保留完整的推导条件及数学论证。这里保留完整的推导条件及数学论证。\[a_{5}=a_{4}+1\]

第6步，这里保留完整的推导条件及数学论证。这里保留完整的推导条件及数学论证。这里保留完整的推导条件及数学论证。\[a_{6}=a_{5}+1\]


% SOURCE-END U0001
}
\end{frame}
% FRAME-ID F004
\begin{frame}[t,fragile]{归纳证明（续2）}
\normalsize
{\normalsize
\textbf{证明：}
% SOURCE-BEGIN U0001 444 590
第7步，这里保留完整的推导条件及数学论证。这里保留完整的推导条件及数学论证。这里保留完整的推导条件及数学论证。\[a_{7}=a_{6}+1\]

第8步，这里保留完整的推导条件及数学论证。这里保留完整的推导条件及数学论证。这里保留完整的推导条件及数学论证。\[a_{8}=a_{7}+1\]
% SOURCE-END U0001
}
\end{frame}
\end{document}
